% ============================================
% DOC-17 (ZH)
% 中文翻译 — SPEC-935-R356
% ============================================
\section*{\textbf{DOC-17} --- 一个士兵的信}
\markboth{DOC-17 --- 一个士兵的信}{}

\noindent \textit{1915 —— 从里面看到的牲畜圈}

\noindent \textbf{虚构重构 —— 叙事文件}
\textbf{1915 —— 无已知档案原件的虚构书信}

\vspace{0.5em}\noindent\rule{\textwidth}{0.5pt}\vspace{0.5em}

\noindent \textit{虚构书信，署名安东尼奥·阿尔维斯·内托士兵，设定为驻守福塔莱萨附近的阿拉加迪索（Alagadiço）收容营。}

\noindent \textit{虚构重构，寄给他在索布拉尔的妻子玛丽亚；未在军事档案中发现对应原件。}

\vspace{0.5em}\noindent\rule{\textwidth}{0.5pt}\vspace{0.5em}

\noindent \textit{玛丽亚，}

\noindent \textit{我在牲畜圈的帐篷里写这封信，但不知道能不能寄出去。这里的情况比上校向媒体说的严重得多。}

\noindent \textit{每天都有更多的逃荒者到来。整个整个的家庭。走了好几天、没有食物没有水的人。孩子们到达时眼窝深陷，肚子鼓胀。女人们一到就被分开，送去“重新安置”。我知道那意味着什么。所有人都知道。}

\noindent \textit{男人和孩子留下。留在铁丝网后面，等待永远不会来的食物。}

\noindent \textit{政府送牛来，但牛进了我们的厨房，不是他们的。他们每天领到一勺稀汤。这汤不足以让他们一下子死掉。却足以让他们慢慢死掉。}

\noindent \textit{最糟的是孩子们。这里有个男孩，最多十岁，独自一人来的。母亲被卡车带走了。父亲死在路上。他不哭。他只是看着。他的眼睛是黄的。右边那只。他说是“一个怪物”改变的。他不哭，玛丽亚。他只是看着。而那目光……穿透了我。}

\noindent \textit{上周，我在乱葬坑看见了一个东西。}

\noindent \textit{牲畜圈后面埋尸体的坑。我值夜班时听见一个声音。咀嚼声。我去查看。}

\noindent \textit{里面有个东西。住在尸体中间。很大。非常大。有在黑暗中发光的黄眼睛。它没有攻击我。它只是看着。然后继续吃。}

\noindent \textit{我报告了中士。他说那是“幻觉”和“饥饿”。说饥饿会让人看见不存在的东西。}

\noindent \textit{也许吧。也许不是。}

\noindent \textit{但今晚，当我闭上眼睛，我会看见那双黄眼睛。}

\noindent \textit{我看见了。}

\noindent \textit{而且，玛丽亚，它们也看见了我。}

\noindent \textit{--- 安东尼奥}

\vspace{0.5em}
\noindent\textit{\textbf{↪ 路线:} \rota{DOC-18} \rotasep \rota{MEM-06}}
